% Generated draft: requires mathematical and source review.
\documentclass[11pt]{article}
\usepackage[margin=25mm]{geometry}
\usepackage{fontspec}
\setmainfont{Latin Modern Roman}
\usepackage{amsmath,amssymb,mathtools,unicode-math}
\setmathfont{Latin Modern Math}
\usepackage{enumitem,needspace,xurl,hyperref}
\setlength{\parindent}{0pt}
\setlength{\parskip}{5pt}
\setlength{\emergencystretch}{3em}
\newcommand{\N}{\mathbb{N}}
\newcommand{\Z}{\mathbb{Z}}
\newcommand{\Q}{\mathbb{Q}}
\newcommand{\R}{\mathbb{R}}
\newcommand{\C}{\mathbb{C}}
\newcommand{\F}{\mathbb{F}}
\newcommand{\A}{\mathbb{A}}
\newcommand{\PP}{\mathbb{P}}
\newcommand{\E}{\mathbb{E}}
\newcommand{\eps}{\varepsilon}
\newcommand{\dd}{\,\mathrm{d}}
\begin{document}
\section*{UnsolvedMath ID: 30005658}
\textbf{Source problem code:} catalogue entry
\subsection{Definitions and mathematical statement}
Let \(A\) be a connective bounded commutative differential graded ring, and let \(\mathcal D=D(A)\) be its unbounded derived category of differential graded modules, equipped with derived tensor product \(\otimes_A^{\mathbf L}\) and tensor unit \(\mathbf 1=A\). With cohomological grading, the usual connective convention is \(H^i(A)=0\) for \(i>0\); cohomological boundedness means that \(H^i(A)=0\) outside a finite interval. The supplied record does not explicitly specify these grading and boundedness conventions. Put
\[
R=\operatorname{End}_{\mathcal D}(\mathbf 1)\cong H^0(A),\qquad X=\operatorname{Spec}(R).
\]
Here \(\operatorname{End}\) denotes degree-zero endomorphisms, and \(X\) has its Zariski topology.

A \(t\)-structure is a pair \((\mathcal U,\mathcal V)=(\mathcal D^{\leq 0},\mathcal D^{\geq 0})\) of full subcategories such that \(\mathcal U[1]\subseteq\mathcal U\), \(\mathcal V[-1]\subseteq\mathcal V\), \(\operatorname{Hom}_{\mathcal D}(\mathcal U,\mathcal V[-1])=0\), and every \(M\in\mathcal D\) belongs to a distinguished triangle
\[
U\longrightarrow M\longrightarrow W\longrightarrow U[1],
\qquad U\in\mathcal U,\quad W\in\mathcal V[-1].
\]
An object \(C\) is compact if \(\operatorname{Hom}_{\mathcal D}(C,-)\) commutes with arbitrary coproducts. The \(t\)-structure is compactly generated if there is a set \(\mathcal S\) of compact objects such that
\[
\mathcal V[-1]=\bigl\{Y\in\mathcal D:\operatorname{Hom}_{\mathcal D}(C[k],Y)=0\text{ for every }C\in\mathcal S\text{ and every }k\geq 0\bigr\}.
\]

\textbf{Source question.} For every connective bounded commutative differential graded ring \(A\), are the compactly generated \(t\)-structures on \(D(A)\) whose nonpositive aisle is a \(t\)-\(\otimes\)-ideal classified by filtrations by supports
\[
f:\mathbb Z\longrightarrow 2^X,\qquad X=\operatorname{Spec}\bigl(\operatorname{End}_{D(A)}(\mathbf 1)\bigr)?
\]
Here \(2^X\) is the power set of \(X\). No nondegeneracy, boundedness, or completeness condition on the \(t\)-structures is supplied.

\textbf{Unresolved qualifications.} The original statement uses the term \(t\)-\(\otimes\)-ideal, whereas the clean statement says tensor ideal. These do not specify the same condition without further conventions. Closure under \(U\otimes_A^{\mathbf L}M\) for every \(U\in\mathcal U\) and every \(M\in D(A)\) is a stronger requirement than closure only for \(M\) in a specified nonpositive aisle; the latter also requires a reference \(t\)-structure to be identified. The supplied record does not settle the intended condition.

The record also does not define a filtration by supports, its monotonicity, or the support construction giving the proposed classification. A customary possible interpretation is a decreasing sequence \(f(n+1)\subseteq f(n)\) of Thomason subsets, where a Thomason subset is a union of closed subsets with quasi-compact open complements. Another possible support condition is specialization closure. These possibilities are recorded for clarification and are not adopted as missing hypotheses. No endpoint or exhaustiveness condition on \(f\) is supplied. Thus the displayed question preserves the supplied target, but its precise parameter space and classification correspondence require clarification before it becomes a fully specified bijection statement.
\subsection{Short English statement}
For a connective bounded commutative differential graded ring \(A\), the question asks whether support filtrations on \(\operatorname{Spec}(H^0(A))\) classify compactly generated \(t\)-structures on \(D(A)\) satisfying the intended tensor compatibility condition. The meaning of that condition, the permitted support filtrations, and the classification correspondence remain unspecified. The supplied record labels the problem open; this draft does not independently verify that label.
\subsection{Sources}
\begin{itemize}
\item Local UnsolvedMath registry: \texttt{data/unsolvedmath/problems.json}, permanent ID \(30005658\), title \emph{Support-Filtration Classification of Compactly Generated t-Structures}, exact source problem code \texttt{catalogue entry}. The supplied record gives proposed year \(2023\) and names no individual proposer.
\item Supplied report citation: \emph{Tensor-Triangular Geometry and Interactions (2024)}, Oberwolfach Reports, DOI \texttt{10.4171/owr/2023/40}, pp.\ 2303--2358. \url{https://doi.org/10.4171/owr/2023/40}. The supplied background attributes the question to workshop open-problem discussions documented in Oberwolfach Reports. The citation's year \(2024\) and the record's proposed year \(2023\) are retained without reconciliation.
\item Supplied literature source: Hrbek, \emph{Compactly generated t-structures}. \url{https://arxiv.org/abs/1806.00078}.
\item Supplied literature source: OWR report. \url{https://publications.mfo.de/bitstream/handle/mfo/4102/OWR_2023_40.pdf?sequence=4}.
\item Literature assessment dated \texttt{2026-08-22}: the supplied curation labels the problem open and reports that support-filtration classifications are established for ordinary commutative rings, while no general theorem was found for the stated connective bounded commutative differential graded setting. This is supplied dated curation derived from the OWR material, not an independently verified literature assessment; its scope and time-sensitive claims require independent checking.
\item The registry labels the statement \texttt{corrected\_verified} and reports that the official report was checked directly and that the clean statement restores definitions while preserving the mathematical target. These are supplied source-verification claims, not verification performed for this draft.
\end{itemize}
\end{document}
